医学图像分割是辅助疾病诊断与治疗规划的关键环节，但其深度学习方法的发展长期受限于高质量标注数据的匮乏。临床标注过程不仅耗时费力，不同医疗机构之间的数据分布也差异显著，这使得现有模型在跨中心、跨模态场景下泛化性能较差，难以真正落地应用。

针对标注稀缺与模型泛化能力不足这两个核心问题，本研究系统构建了一套自适应的医学图像分割理论与方法体系。一方面，为提升有限标注数据的利用效率，设计了认知不确定性引导的半监督学习框架，通过量化模型预测的不确定性来优化伪标签质量；另一方面，为增强模型对多尺度特征的提取能力，深入探究了U-Net与Transformer的融合范式，并在网络的跳跃连接中引入通道注意力机制，以强化关键特征的传递。

为验证所提方法的有效性，本研究基于CA数据集与Synapse数据集开展了系列对比实验。在半监督实验部分，所提UA-MT框架仅使用20\%的有标签数据（16例）即可达到88.18\%的Dice系数，显著优于标准Mean Teacher方法。在网络架构实验部分，引入双重注意力机制的DA-TransUNet模型在多器官分割任务中取得了80.67\%的平均Dice系数，且边界分割精度（HD95）明显优于U-Net及Swin-Unet等基准模型。上述结果表明，本研究提出的方法能够有效缓解标注数据稀缺带来的瓶颈，并显著提升分割精度与模型鲁棒性。